% GDZ PPN236018647,native522lower–524;HKUST523–525.
\NumberedItem{8}
\indent Scalam simplicem in seriebus variatim recurrentibus esse functionem similem secundi ordinis scalarum componentium.
\par\noindent\hfill[1796] 26. Mai.

\medskip\begingroup\small
\indent Bedeutet \(G(x)\) eine ganze rationale Funktion von \(x\) von \((n-1)\)-tem oder niedrigerem Grade und entwickelt man
\[
\frac{G(x)}{1-a_1x-a_2x^2-\ldotp\,\ldotp\,\ldotp-a_nx^n}
\]
\noindent in die unendliche Reihe \(s_0+s_1x+s_2x^2+\ldotp\,\ldotp\,\ldotp\), so ist
\[
s_{n+t}=a_1s_{n+t-1}+a_2s_{n+t-2}+\ldotp\,\ldotp\,\ldotp+a_ns_t\qquad(t=0,1,2,\ldotp\,\ldotp\,\ldotp).
\]
\noindent Eine solche Reihe nennt A. \textsc{de Moivre} (\textit{Miscellanea analytica de seriebus et quadraturis}, Londini 1730, S. 27, siehe auch schon Philosophical Transactions, London 1722, S. 162) \textls[120]{rekurrent} und bezeichnet die Zahlen \(a_1,a_2,\ldotp\,\ldotp\,\ldotp,a_n\) als \textls[120]{Index} oder \textls[120]{Scala relationis}\begingroup\renewcommand{\thefootnote}{***)}\footnote{Vergl. auch die Kapitel IV, XIII und XVII von \textsc{Eulers} \textit{Introductio in analysin infinitorum} I, Lausannae 1748, wo \textsc{Moivre} angeführt und die Bezeichnung scala gebraucht wird, sowie desselben Verfassers \textit{Institutiones calculi differentialis} I, 1755, § 46, L. Euleri Opera omnia, ser. I, vol. 10, S. 45.}\endgroup. Sowohl hier als auch in den Aufzeichnungen Nr. 10 und Nr. 20, die sich auf denselben Gegenstand beziehen, denkt sich \textsc{Gauss} anscheinend den Nenner in Faktoren zerlegt und die zum Nenner gehörige \textls[120]{einfache} Skala aus den zu den Faktoren gehörigen Skalen \textls[120]{komponiert}. Hier in der Nr. 8, wo von einer Funktion zweiten Grades die Rede ist, hat er vermutlich die Zerlegung des Nenners in zwei Faktoren im Auge, wogegen sich die Nr. 20 auf den Fall von drei und mehr Faktoren beziehen könnte. Die Zerlegung des Nenners in ein Produkt von Faktoren ersten Grades und die Entwicklung der reziproken Werte dieser Faktoren in geometrische Reihen kommt im art. 1. der aus dem Nachlaß herausgegebenen \textit{Theoria interpolationis methodo nova tractata} (Werke III, S. 265) vor. Ein Beispiel einer rekurrenten Reihe für \(n=2,\ a_1=\frac12,\ a_2=-\frac12\) findet sich im art. 4. des nachgelassenen Bruchstücks \textit{Grundbegriffe der Lehre von den Reihen} (oben S. 393).
\par\noindent\hfill\textsc{Klein. Loewy.}
\par\endgroup

\NumberedItem{9}
\indent Comparationes infinitorum in numeris primis et factoribus cont[entorum].
\par\noindent\hfill[1796] 31. Mai.\quad G[ottingae]

\medskip\begingroup\small
\indent Diese Aufzeichnung, sowie die Nr. 13 vom 19. Juni 1796 beziehen sich wohl auf die oben S. 11, 12 abgedruckten Eintragungen in \textsc{Schulzes} \textit{Tafeln}, die vom Mai 1796 datiert sind.
\par\noindent\hfill\textsc{Bachmann}.
\par\endgroup

\clearpage\SourceLeaf{523}\ArithmeticHead{494}{L}{TAGEBUCH MIT ERLÄUTERUNGEN.}
\NumberedItem{10}
\indent Scala ubi seriei termini sunt producta vel adeo functiones quaecunque terminorum quotcunque serierum.
\par\noindent\hfill[1796] 3. Iun.\quad G[ottingae]
\par\medskip\begingroup\small
\indent Vergl. die Bemerkung zu der Nr. 8.
\par\endgroup

\NumberedItem{11}
\indent Formula pro summa factorum numeri cuiusvis compositi:
\[
\text{f[actum] gener[ale]}\ \frac{a^{n+1}-1}{a-1}.
\]
\par\noindent\hfill[1796] 5. Iun.\quad G[ottingae]

\medskip\begingroup\small
\indent Diese Notiz spricht den Satz aus, daß wenn eine Zahl \(N\) in Primzahlpotenzen zerlegt gleich \(\Pi a^n\) ist, die Summe der Teiler von \(N\) durch das Produkt (factum)
\[
\Pi\frac{a^{n+1}-1}{a-1}
\]
\noindent dargestellt wird.
\par\noindent\hfill\textsc{Bachmann}.
\par\endgroup

\NumberedItem{12}
\indent Periodorum summa omnibus infra modulum numeris pro elementis sumtis:
\[
\text{fact[um] gen[erale]}\ ((n+1)a-na)a^{n-1}
\]
\par\noindent\hfill[1796] 5. Iun.\quad G[ottingae]

\medskip\begingroup\small
\indent Der Ausdruck Periode hat hier sicher die gleiche Bedeutung wie in den \textit{Disquisitiones arithmeticae} art. 46, Werke I, S. 39. Ist \(N\) teilerfremd zu \(p\) und \(d\) der Exponent, zu dem \(N\) modulo \(p\) gehört oder allgemeiner, für den
\[
N^d\equiv1\ (\mathrm{mod.}\,p)
\]
\noindent wird, so nennt dort \textsc{Gauss} die Potenzen
\[
1,\ N,\ N^2,\ \ldotp\,\ldotp\,\ldotp,\ N^{d-1}
\]
\noindent oder ihre Reste modulo \(p\) eine Periode. Die Wendung »omnibus infra modulum numeris pro elementis sumtis« zeigt, daß in unserer Aufzeichnung der Modulus \(p\) als Primzahl gedacht ist. Nun ist für die \textls[120]{Summe} aller dieser Perioden, auch wenn ihre einzelnen Glieder durch die Reste nach dem Modul \(p\) ersetzt werden, ein dem angedeuteten Produkte ähnlicher Ausdruck nicht vorhanden. Würde unter \textls[120]{Periode} nur die Anzahl ihrer Glieder verstanden, so wäre die periodorum summa, da \(\varphi(d)\) Zahlen zum Exponenten \(d\) gehören, die auf alle Teiler \(d\) von \(p-1\) bezogene Summe \(\sum d.\varphi(d)\), für die man, wenn man in Primfaktoren zerlegt
\[
p-1=\Pi a^n
\]
\endgroup
\clearpage\SourceLeaf{524}\ArithmeticHead{495}{R}{1796 IUN. 3.—1796 IUN. 20.}
\begingroup\small
\noindent gedacht wird, den Ausdruck
\[
\Pi\frac{a^{2n+1}+1}{a+1}
\]
\noindent erhält. Zieht man nur die voneinander verschiedenen Perioden in Betracht, so ergibt sich (siehe die vorhergehende Nr. 11)
\[
\sum d=\Pi\frac{a^{n+1}-1}{a-1};
\]
\noindent beides stimmt nicht. Nimmt man aber den Ausdruck \textls[120]{Summa} im Sinne von \textls[120]{Gesamtheit} oder \textls[120]{Anzahl}, so besagte unsere Notiz nichts anderes als den Satz (\textit{Disquis. arithm.} art. 39., Werke I, S. 31)
\[
\sum\varphi(d)=p-1=\Pi a^n,
\]
\noindent was, wenn die Formel von \textsc{Gauss} richtig geschrieben ist, mit ihr übereinstimmen würde; doch wäre dann die dem Faktor des Produkts gegebene Form sehr wunderlich.

\indent Wahrscheinlich liegt aber hier ein Schreibfehler vor und der Faktor soll
\[
((n+1)a-n)a^{n-1}
\]
\noindent lauten. Bildet man nämlich für alle Reste modulo \(p\) oder, was auf dasselbe hinauskommt, für alle Exponenten \(k=1,2,\ldotp\,\ldotp\,\ldotp,p-1\) und für eine primitive Wurzel \(g\) von \(p\) die Potenzen
\[
1,\ g^k,\ g^{2k},\ \ldotp\,\ldotp\,\ldotp,\ g^{(p-2)k},
\]
\noindent so erscheint für die \(\varphi(d)\) Werte von \(k\), die teilerfremd sind zu dem Teiler \(d\) von \(p-1\), je eine \(d\)-gliedrige Periode \(\delta=\dfrac{p-1}{d}\)mal, und man erhält so eine \textls[120]{Summa} — im Sinne von Gesamtheit oder Anzahl — von
\[
\sum_{d\delta=p-1}\varphi(d).\delta
\]
\noindent Perioden. Hier ist aber
\[
\begin{aligned}
\sum_{d\delta=p-1}\varphi(d).\delta
&=(p-1)\Pi\left(1+\frac{a-1}{a}+\frac{a(a-1)}{a^2}+\ldotp\,\ldotp\,\ldotp+\frac{a^{n-1}(a-1)}{a^n}\right)\\
&=\Pi a^n.\Pi\left(1+n\frac{a-1}{a}\right)\\
&=\Pi((n+1)a-n)a^{n-1}.
\end{aligned}
\]
\par\noindent\hfill\textsc{Bachmann}.
\par\endgroup

\NumberedItem{13}
\indent Leges distributionis.
\par\noindent\hfill[1796] 19. Iun.\quad G[ottingae]
\par\medskip\begingroup\small
\indent Vergl. die Bemerkung zu der Nr. 9.
\par\endgroup

\NumberedItem{14}
\indent Factorum Summae in Infinito \(=\dfrac{\pi\pi}{6}.\)Sum[mam] Num[erorum].
\par\noindent\hfill[1796] 20. Iun.\quad G[ottingae]
\par\medskip\begingroup\small
\indent Vergl. die Notiz Nr. 31 vom 6. September 1796 und die oben S. 14 abgedruckte Aufzeichnung aus den \textit{Exercitationes Mathematicae} sowie die zugehörigen Bemerkungen oben S. 17, 18.
\par\noindent\hfill\textsc{Klein. Bachmann.}
\par\endgroup
